This work investigates whether the 24-hour Holter electrocardiogram (ECG)
can identify, among patients with chronic heart failure, those who will
suffer sudden cardiac death (SCD) within four years. The problem was framed
as a patient-level binary classification: from the Holter recording,
distinguish the 55 patients who died of SCD from the 462 who survived the
four-year follow-up, in a cohort of 517 patients in sinus rhythm from the
public MUSIC database (PhysioNet). In stratified five-fold cross-validation with
patient-level splits, the main metric is the
area under the ROC curve (AUC) for predicting SCD versus survival, computed
on out-of-fold predictions, with 95\,\% bootstrap confidence intervals (CI).
Four approaches were compared. (i) Convolutional and residual networks,
trained directly on the three-lead signal to predict SCD, in 19
configurations, reached AUC
between 0.48 and 0.56, not different from chance. (ii) Markers extracted from
the signal --- heart rate variability, ectopy and beat morphology ---, in
logistic regression, predicted SCD with AUC between 0.58 and 0.62 with 44
markers and 0.647 (CI: 0.56--0.73) with two markers defined from the
literature before the analysis: the variation of heart rate over the day and
the ectopic burden. (iii) Six clinical variables
(age, sex, NYHA class, Q waves and two arrhythmias) predicted SCD
with AUC 0.661 (CI: 0.58--0.74). (iv) In the proposed approach, a ResNet-18
was trained to recognize, in 30-second Holter windows, 45 clinical
attributes measured by other exams, such as QRS duration, ejection fraction
and Pro-BNP, and the representation it learned from the signal was used by
a logistic regression to predict SCD. The network learned to recognize
electrical attributes in the Holter, such as QRS above 120~ms (AUC 0.92).
Using only the Holter signal, its representation predicted SCD with AUC
0.682 (CI: 0.60--0.76), equivalent to the clinical variables and 0.17 above
the same network trained directly for SCD (0.508); combined with the
clinical variables, it reached 0.702 (CI: 0.62--0.78). The best model of
this work, with the clinical variables and the measured exams
(echocardiography, Pro-BNP and resting ECG), reached 0.733 (CI: 0.66--0.80),
and no Holter-based approach added information to it. The results show that
guiding the network with clinical attributes is what makes it possible to
extract from the Holter, with no clinical variable as input, an SCD risk
estimate at the level of the clinical variables, which the same network
trained directly for SCD did not achieve.

\textbf{Keywords:} Sudden Cardiac Death. Holter ECG. Deep Learning. Heart Rate
Variability. Heart Failure. Risk Stratification.
